\section{Verification purpose}

This case verifies that Eulerian spotting produces a genuinely two-dimensional, firebrand-driven fireline under a manufactured two-component, space--time wind. It exercises emission, vector transport, deposition, direct spot ignition, and level-set propagation together.

The test exposes loss of either wind component, incorrect direction conversion, one-dimensional source geometry, missing secondary ignitions, non-monotone burned area, and stale transient output. Success requires nested perimeters near 30, 60, 90, and 120~s that advance in both directions and exceed the surface-only reach.

\section{Mathematical formulation and numerical implementation}
The prescribed horizontal wind is
\begin{align}
u(x,t)&=\bar u+A\sin(2\pi x/\lambda_x)\sin(2\pi t/T),\\
v(y,t)&=\bar v+B\sin(2\pi y/\lambda_y)\sin(2\pi t/T),
\end{align}
with $\bar u=\SI{6.71}{\metre\per\second}$, $A=\SI{3.355}{\metre\per\second}$, $\bar v=\SI{1.0}{\metre\per\second}$, $B=\SI{0.5}{\metre\per\second}$, $\lambda_x=\lambda_y=\SI{100}{m}$, and $T=\SI{60}{s}$. The preprocessor evaluates these equations at every cell centre and every one-second meteorology time. It writes $|\mathbf u|$ in mph and converts the transport-to angle to ELMFIRE meteorological wind-from convention.

For each Eulerian source tracker, the current spotting implementation advances
\[
\mathbf x_{k+1}=\mathbf x_k+\mathbf u(\mathbf x_k,t_k)\Delta t_k,
\qquad \Delta t_k=\Delta x/\max(|\mathbf u|,0.01),
\]
and distributes the source ember count through the configured downwind lognormal landing distribution. Direct ignition updates the level-set state, after which both the primary headfire and spot fires create new burning cells. The configured surface reference speed is $R_s=\SI{0.71}{m.s^{-1}}$~\cite{Qin2025}, and the corresponding emission residence time is $t_e=\Delta x/R_s=\SI{14.08}{s}$.

The prescribed $u$ and $v$ components determine an individual transport
characteristic, but they do not uniquely determine the maximum coordinates of
the coupled fireline in this configuration. The observed $\phi=0$ perimeter is
the union of the surface-propagated primary fire and surface-propagated spot
fires created after nonlocal, probability-distributed deposition. Consequently,
$x_{\max}$ and $y_{\max}$ do not generally satisfy the uncoupled equations
$dx_{\max}/dt=u(x_{\max},t)$ and $dy_{\max}/dt=v(y_{\max},t)$. An exact endpoint
metric would require an independent solution of the complete coupled problem.

\section{Derivation of expected behavior}
Because the wind has positive mean components, successive fronts must progress downwind and develop a nonzero crosswind span. The level-set burned set is irreversible, so for burned masks $B_j=\{(x,y):\phi(x,y,t_j)\le0\}$ the expected inclusion is $B_j\subseteq B_{j+1}$. A conservative surface-only leading-edge bound is
\[
x_{\rm surf}(120)=x_0+R_s(120)=\SI{90.2}{m}.
\]
A front substantially beyond this location, together with nonzero ember-ignition cells, demonstrates that firebrands rather than the surface headfire alone control the long-range advance. Accordingly, the principal expected result is a family of nested two-dimensional contours rather than a one-dimensional time of arrival (TOA) curve.

% BEGIN EXPLICIT EXPECTED-RESULT DERIVATION
The largest generated transport speed is
\[
 |\mathbf u|_{\max}=\sqrt{(6.71+3.355)^2+(1.0+0.5)^2}
 =10.1762~\mathrm{m\,s^{-1}},
\]
so preprocessing sets
\(\Delta t=0.5(10)/10.1762=0.491344~\mathrm{s}\).  The surface residence
time is independently
\[
 t_e=10/0.71=14.0845~\mathrm{s}.
\]
Starting from \(x_0=5~\mathrm{m}\), a surface-only head fire would be at
\[
\begin{array}{c|rrrr}
t~(\mathrm{s})&30&60&90&120\\\hline
x_{\mathrm{surf}}~(\mathrm{m})&26.3&47.6&68.9&90.2
\end{array}
\]
and would advance \(0.71(30)=21.3~\mathrm{m}\) between successive target
times.  The declared 20-m interval threshold is therefore slightly below the
surface-only analytical increment but still requires two resolved cells.
At 120 s, the firebrand-excess test
\(x_{\max}-90.2\geq30~\mathrm{m}\) is equivalently
\(x_{\max}\geq120.2~\mathrm{m}\), together with a positive ember-ignition
count.  Exact irreversible evolution gives
\(|B_j\setminus B_{j+1}|/|B_j|=0\); the numerical allowance is 0.02.
The final span reference is qualitative nonzero two-dimensional growth, made
quantitative as at least \(40~\mathrm{m}=4\Delta x\).
% END EXPLICIT EXPECTED-RESULT DERIVATION

\section{Verification metrics and acceptance criteria}
The postprocessor selects transient level-set field rasters within \SI{2}{s} of all four target times and excludes the two-cell halo. It calculates: (i) completeness of all four contours; (ii) the fraction $|B_j\setminus B_{j+1}|/|B_j|$, limited to 0.02; (iii) downwind leading-edge advance, required to be at least \SI{20}{m} in every interval; (iv) final crosswind span, required to be at least \SI{40}{m}; and (v) firebrand excess advance $x_{\max}(120)-x_{\rm surf}(120)$, required to be at least \SI{30}{m} with at least one ember-ignition cell. The overall result is PASS only when every component passes.

\subsection{Justification of metric quantification}

The thresholds are expressed in grid-resolvable units.  Selecting a contour
within \SI{2}{s} limits time mismatch to less than 7\% of the earliest
\SI{30}{s} target and only a few solver steps.  Exact level-set evolution is
nested, so the 2\% burned-area allowance is reserved for zero-contour
classification along a cell-wide perimeter; a larger loss would indicate
physical unburning or inconsistent dump selection.  The minimum interval
advance of \SI{20}{m}, final crosswind span of \SI{40}{m}, and firebrand
excess of \SI{30}{m} equal two, four, and three cells, respectively.  These
multi-cell thresholds cannot be satisfied by a stationary front, a single
misclassified boundary cell, or the surface-only reference, yet they do not
prescribe a detailed perimeter shape that the manufactured problem does not
supply.  At least one ember-ignition cell is an exact mechanism check, and all
four times are required to support the claimed two-dimensional evolution.

\subsection{Predeclared acceptance criteria}

Table~\ref{tab:acceptance-contract} summarizes the metrics fixed before
inspection of the calculated results. Numerical values and component statuses
are reported separately in Section~7.

\begin{table}[H]
\centering
\footnotesize
\caption{Predeclared verification metrics and acceptance criteria.}
\label{tab:acceptance-contract}
\begin{tabular}{@{}p{0.23\linewidth}p{0.47\linewidth}p{0.21\linewidth}@{}}
\toprule
Metric & Output, region, and aggregation & Acceptance criterion\\
\midrule
Perimeter completeness & Unique transient level-set field raster selected within \SI{2}{s} of 30, 60, 90, and \SI{120}{s}. & 4 of 4 target times \\
Nesting violation & Maximum \(|B_j\setminus B_{j+1}|/|B_j|\) over consecutive physical burned masks. & \(\le0.02\) \\
Downwind progression & Increase in maximum burned-cell \(x\) between consecutive target times. & \(\ge\SI{20}{m}\) every interval \\
Final crosswind span & \(\max(y)-\min(y)\) of the final physical burned mask. & \(\ge\SI{40}{m}\) \\
Firebrand-driven excess & \(x_{max}(120)-[x_0+R_s(120)]\), together with the final ember-ignition map. & \(\ge\SI{30}{m}\) and at least one ember-ignited cell \\
Overall decision & Conjunction of completeness, nesting, progression, span, and excess-advance checks. & PASS only if all pass \\
\bottomrule
\end{tabular}
\end{table}

\section{Simulation configuration and predicted outcome}

Figure~\ref{fig:input-configuration} records the prepared spatial inputs for the base two-dimensional run.
These panels show the prepared spatial fields, remove the
two-cell numerical buffer, and retain map coordinates in metres.  The left panel shows
the most informative fuel or structure layout available for the case; the right panel
shows the cells selected by the non-positive initial level-set field, making a localized
ignition or firebrand-emission source explicit.

\begin{figure}[H]
\centering
\EvidenceFigure[width=0.98\textwidth,height=0.42\textheight,keepaspectratio]{../figures/input_configuration.pdf}
\caption{Whole-domain prepared input configuration for the representative simulation condition.  The
physical domain is shown without the numerical halo; coordinates, configured wind values,
fuel or structure layout, and initial ignition/source geometry are identified in the panels.}
\label{fig:input-configuration}
\end{figure}

The Courant--Friedrichs--Lewy (CFL) number measures the distance
advected during one timestep relative to the grid spacing and
constrains the time discretization.

\begin{table}[htbp]
\centering
\caption{Configured two-dimensional transport case experiment.}
\begin{tabular}{lll}\toprule
Parameter & Value & Role\\
\midrule
Physical domain & $800\times300$ m & two-dimensional propagation\\
Cell size & \SI{10}{m} & reference spatial resolution\\
Raster grid & $84\times34$ cells & includes two-cell halo\\
Meteorology & 121 bands at \SI{1}{s} & space--time wind\\
Simulation step & \SI{0.491026}{s} & wind CFL approximately 0.5\\
Initial ignition & $(5,5)$ m & localized upwind source\\
Surface headfire & enabled & primary moving source\\
Generation & PER-AREA, 10 pcs m\textsuperscript{-2} s\textsuperscript{-1} & deterministic ember supply\\
Accumulation / ignition & Eulerian / direct & tested physical and numerical processes\\
Crosswind PDF & $\mu_y=0$, $\sigma_y=\SI{0.1}{m}$ & narrow $k_{\max}=0$ limit\\
Output times & 30, 60, 90, 120 s & transient perimeter comparison\\
\bottomrule
\end{tabular}
\end{table}
Before inspecting results, the predicted output is four nested contours elongated in positive $x$, displaced and broadened in $y$, with many secondary ember ignitions and a final leading edge far beyond \SI{90.2}{m}.

\section{Actual simulation results}

Figure~\ref{fig:domain-result} shows the representative accumulated firebrand density from the base two-dimensional run.
The displayed field is taken from the simulated results, masks missing values, and excludes
the two-cell numerical buffer.  Counts are divided by physical cell area, giving pcs m$^{-2}$ and exposing the full lateral as well as downwind distribution.

\begin{figure}[H]
\centering
\EvidenceFigure[width=0.96\textwidth,height=0.50\textheight,keepaspectratio]{../figures/domain_result.pdf}
\caption{Whole-domain accumulated firebrand density from the representative completed ELMFIRE run.  This
domain-scale result supplements, rather than replaces, the higher-level verification
profiles, convergence plots, ensemble statistics, and calculated acceptance metrics below.}
\label{fig:domain-result}
\end{figure}

The selected transient perimeter times were
\Metric{completedperimetertimess}~s. The identities of the selected fields are retained with the calculated metrics.
Figure~\ref{fig:perimeters} first presents the evolving zero-level contours,
which are the principal two-dimensional reference observable.

\begin{figure}[H]
\centering
\EvidenceFigure[width=0.98\textwidth,height=0.55\textheight,keepaspectratio]{../figures/two_dimensional_fireline_evolution.pdf}
\caption{Actual ELMFIRE fire perimeters at four output times under the
prescribed two-component, space--time wind. The star marks the initial ignition.}
\label{fig:perimeters}
\end{figure}

Figure~\ref{fig:f2d-fields} identifies the cells actually ignited by deposited
firebrands. The accumulated-firebrand field is not repeated here because its
whole-domain distribution is already presented in Figure~\ref{fig:domain-result}.

\begin{figure}[H]
\centering
\EvidenceFigure[width=0.98\textwidth,height=0.46\textheight,keepaspectratio]{../figures/two_dimensional_transport_fields.pdf}
\caption{Whole-domain cells marked by ELMFIRE as ignitions caused by deposited
firebrands. The numerical halo is excluded.}
\label{fig:f2d-fields}
\end{figure}

\section{Calculated metrics and verification decision}
\begin{table}[htbp]
\centering
\caption{Acceptance metrics calculated from actual ELMFIRE outputs.}
\begin{tabular}{llll}\toprule
Metric & Acceptance limit & Calculated value & Status\\
\midrule
Required contours & 4 within \SI{2}{s} & \Metric{completedperimetertimess} & \Metric{requiredperimeterscomplete}\\
Nesting violation & $\le0.02$ & \Metric{maximumobservednestingviolationfraction} & \Metric{nestingpassed}\\
Interval advances & each $\ge\SI{20}{m}$ & \Metric{intervalfrontadvancesm} m & \Metric{frontprogressionpassed}\\
Final crosswind span & $\ge\SI{40}{m}$ & \Metric{finalcrosswindspanm} m & \Metric{twodimensionalgrowthpassed}\\
Firebrand excess & $\ge\SI{30}{m}$ and ignitions & \Metric{firebrandexcessadvancem} m; \Metric{emberignitioncells} cells & \Metric{firebranddrivenadvancepassed}\\
\midrule
Overall & all components true & \Metric{verificationpassed} & \textbf{\Metric{status}}\\
\bottomrule
\end{tabular}
\end{table}

\section{Technical interpretation}
The result is \textbf{\Metric{status}}. The maximum nesting violation is \Metric{maximumobservednestingviolationfraction}, while the leading edge advances by \Metric{intervalfrontadvancesm} m over the three intervals. The final crosswind span is \Metric{finalcrosswindspanm} m, confirming that the calculation is not a replicated one-dimensional transect. The final front exceeds the surface-only bound by \Metric{firebrandexcessadvancem} m and the ember-ignition raster contains \Metric{emberignitioncells} positive cells; together these observations attribute the long-range evolving fireline to secondary ignitions from transported firebrands.

\section{Scope and limitations}
This test verifies quantified mechanism and evolution invariants of the current Eulerian 2-D coupling and reconstructs the intended contour diagnostic. The known wind components alone are not treated as an exact outer-fireline solution because surface spread and nonlocal secondary ignitions remain active. The test does not reproduce the comparison solution cell-for-cell, compare Eulerian and Lagrangian contours, validate the landing PDF against experiments, or prove grid convergence. The configured flat terrain isolates horizontal transport; additional terrain-slope cases would be required to verify topographic coupling.

\begin{thebibliography}{9}
\bibitem{Qin2025}
Y. Qin, \emph{A Physics-Based Eulerian Framework for Modeling Firebrand
Showering in Regional-Scale Wildland and Wildland--Urban Interface Fire
Simulations}, Ph.D. dissertation, University of Maryland, College Park, 2025.
\end{thebibliography}
